%% ========================================================================
%% Appendix C: Glossary
%% ========================================================================

\chapter{Glossary}
\label{app:glossary}

This appendix contains definitions of technical terms used in this book.

%% ------------------------------------------------------------------------
\section*{A}

\begin{description}
    \item[ACL (Access Control List)] A more granular access control mechanism
          than traditional Unix permissions, allowing access rights configuration
          for multiple users and groups.

    \item[API (Application Programming Interface)] A set of definitions and protocols
          for building and integrating application software.
\end{description}

%% ------------------------------------------------------------------------
\section*{B}

\begin{description}
    \item[Bash (Bourne Again Shell)] A command-line shell and programming language
          that is the default in most Linux distributions.

    \item[Bootloader] A program that loads the operating system into memory when
          the computer is powered on. Examples: GRUB, LILO.
\end{description}

%% ------------------------------------------------------------------------
\section*{C}

\begin{description}
    \item[CLI (Command Line Interface)] A text-based interface for interacting
          with the operating system using commands.

    \item[Cron] A daemon for scheduling command or script execution at
          specific times.
\end{description}

%% ------------------------------------------------------------------------
\section*{D}

\begin{description}
    \item[Daemon] A program that runs in the background and has no
          direct interface with the user.

    \item[Distro (Distribution)] A Linux variant packaged with a kernel,
          utilities, and specific application software.
\end{description}

%% [Continue for letters E-Z...]

%% ------------------------------------------------------------------------
\section*{K}

\begin{description}
    \item[Kernel] The core of an operating system that manages hardware and
          provides basic services for other software.
\end{description}

%% ------------------------------------------------------------------------
\section*{S}

\begin{description}
    \item[Shell] A program that provides an interface for users to interact
          with the operating system.

    \item[System Call] A programming interface for services provided
          by the operating system kernel.

    \item[Systemd] A modern init system used to bootstrap user space
          and manage system services.
\end{description}

%% [Add other terms...]
